\originalchapter{8.511作业5--11：运动、关联与配对}
\psource{Patrick A. Lee，8.511 Fall 2004，公开Problem Sets5--11。}
\originalsection{作业5（2004年10月19日截止）}
\begin{exercise}[原题1(a)--(c)]
一维周期势的带隙为$\Delta$，在均匀电场$\mathcal E$中价带与导带倾斜，等能的两点A、B可发生Zener隧穿。(a)证明在禁带中仍可作$\Psi_k=\sum_Ga_{k-G}\e^{\ii(k-G)x}$展开而$k$可复。令$G_0=2\pi/a,k=G_0/2+\kappa,E=E_0+\varepsilon,E_0=\hbar^2G_0^2/(8m)$，在两波近似下求$\kappa^2$。(b)若$\varepsilon(x)=e\mathcal E x$，求A、B距离。(c)按WKB求$P=\exp[-2\int_B^A|\kappa(x)|\dd x]$，给出半经典单带图像失效条件。
\end{exercise}
\begin{solution}
(a)Bloch平移关系可解析延拓到复$k$，禁带中复波矢描述衰减而不代表周期盒的可归一本征态。两波矩阵在区界附近对角$E_0\pm\hbar^2G_0\kappa/(2m)$，略去共同$O(\kappa^2)$项，得到$\varepsilon^2=|V_{G_0}|^2+\hbar^2(2E_0/m)\kappa^2$，即$\kappa^2=2m(\varepsilon^2-|V_{G_0}|^2)/(4\hbar^2E_0)$。$|\varepsilon|<|V|$时$\kappa$虚，$\Delta=2|V|$。

(b)$x_\pm=\pm|V|/(e\mathcal E)$，距离$\Delta/(e\mathcal E)$。(c)积分是椭圆面积：$\int_{-|V|}^{|V|}\sqrt{|V|^2-\varepsilon^2}\dd\varepsilon=\pi|V|^2/2$。因此
\begin{equation}
P=\exp\left[-\frac{\pi\sqrt{m/(2E_0)}|V|^2}{\hbar e\mathcal E}\right]=\exp\left[-\frac{\pi^2\Delta^2}{8E_0e\mathcal E a}\right].
\end{equation}
指数不再远大于1时带间隧穿不可忽略。WKB本身要求慢变场与局部带描述，不能把该近似指数当作任意强场的精确概率。
\end{solution}
\begin{exercise}[原题2(a)--(b)]
电子位于边长$L$的立方盒中，磁场$B\hat z$，采用原题Gaussian制$H=(\hbar\nabla/\ii+eA/c)^2/(2m)$与$A=-By\hat x$。(a)求能量。(b)求每级简并，并与Onsager规则比较。
\end{exercise}
\begin{solution}
平面波$\e^{\ii(k_xx+k_zz)}$使$y$方向化为中心$y_0=\hbar ck_x/(eB)$的谐振子，$E_{n,k_z}=\hbar\omega_c(n+1/2)+\hbar^2k_z^2/(2m)$，$\omega_c=eB/(mc)$。$k_x$间距$2\pi/L$，中心落在$0<y_0<L$的数量$eBL^2/(hc)$，不含自旋。固定$k_z$时$A_k=\pi k_\perp^2$，Onsager$A_k=(2\pi eB/\hbar c)(n+1/2)$给完全相同能量和一单位磁通对应一态的计数。
\end{solution}
\originalsection{作业6（2004年10月26日截止，修正版）}
\begin{exercise}[原题1(a)--(b)]
将K近似为自由电子金属。(a)求Shubnikov--de Haas振荡的$1/B$周期。(b)求$B=1\,\mathrm T$时实空间极值轨道面积。原题未给晶格常数或密度，给出依$n$的答案并明确数值假设。
\end{exercise}
\begin{solution}
$A_k=\pi k_F^2$，SI制$\Delta(1/B)=2e/(\hbar k_F^2)$，$A_r=(\hbar/eB)^2A_k$。若取常见室温近似密度$0.862\,\mathrm{g/cm^3}$、摩尔质量$39.0983\,\mathrm{g/mol}$且每原子一价电子，$n=1.328\times10^{28}\,\mathrm{m^{-3}}$，$k_F\simeq7.325\times10^9\,\mathrm{m^{-1}}$，周期约$5.67\times10^{-5}\,\mathrm{T^{-1}}$，$A_r\simeq7.30\times10^{-11}\,\mathrm{m^2}$。这些数字是所列模型输入的计算，不是源题给定数据；实材料参数随温度而变。
\end{solution}
\begin{exercise}[原题2(a)--(d)]
能带$E=\hbar^2[k_x^2/m_x+k_y^2/m_y+k_z^2/m_z]/2$。在半经典方程左侧加$k/\tau$，令$E,B$恒定、$B=B\hat z$、$j=-env$。(a)求$3\times3$DC电导矩阵，以$\omega_{cx}=eB/(m_xc),\omega_{cy}=eB/(m_yc)$表示。(b)求电阻矩阵。(c)讨论低高场Hall电导与Hall电阻。(d)画对角电导和电阻随场变化。
\end{exercise}
\begin{solution}
令$D=1+\omega_{cx}\omega_{cy}\tau^2$。稳态方程为$v_x+\tau\omega_{cx}v_y=-e\tau E_x/m_x$、$v_y-\tau\omega_{cy}v_x=-e\tau E_y/m_y$。解得
\begin{equation}
\sigma=ne^2\tau\begin{pmatrix}(m_xD)^{-1}&-eB\tau/(cm_xm_yD)&0\\eB\tau/(cm_xm_yD)&(m_yD)^{-1}&0\\0&0&m_z^{-1}\end{pmatrix},\quad
\rho=\begin{pmatrix}m_x/(ne^2\tau)&B/(nec)&0\\-B/(nec)&m_y/(ne^2\tau)&0\\0&0&m_z/(ne^2\tau)\end{pmatrix}.
\end{equation}
低场$\sigma_{xy}\propto-B$，高场$\sigma_{xy}\simeq-nec/B$；$\rho_{xy}=B/(nec)$始终线性。$\sigma_{xx},\sigma_{yy}$按$1/(1+\mathrm{const}\,B^2)$下降，$\sigma_{zz}$常数；所有对角$\rho$均常数，该单椭球、常$\tau$模型没有磁电阻，虽有磁电导。SI制把每个$B/c$换成$B$。正文与数值图采用SI，原题表达保留Gaussian换算说明。
\end{solution}
\originalsection{作业7（2004年11月2日截止）}
原题1仅指“Marder第17章题6”，另给$ j_Q=\Pi j-\kappa\partial_xT$、$\Pi_p>0,\Pi_n<0$。完整装置、边界与要求不在公开题面中，故保留引用缺口；不能仅凭热流式猜题。
\begin{exercise}[原题2(a)--(b)]
(a)写N个相互作用电子HF基态能$E_N$，说明为何不等于一电子能级之和。(b)移去第$i$轨道电子并冻结其余轨道，证明$E_N-E_{N-1}(i)=\epsilon_i$，即原题定义的有符号电离能Koopmans关系。
\end{exercise}
\begin{solution}
$E_N=\sum_ih_{ii}+\frac12\sum_{ij}(J_{ij}-K_{ij})$，而$\epsilon_i=h_{ii}+\sum_j(J_{ij}-K_{ij})$。双和每对两次，故$E_N=\sum_i\epsilon_i-\frac12\sum_{ij}(J_{ij}-K_{ij})$。删除$i$后差为$h_{ii}+\sum_j(J_{ij}-K_{ij})=\epsilon_i$，自作用$J_{ii}-K_{ii}=0$。实际移走电子所需的正功常定义为$E_{N-1}-E_N=-\epsilon_i$，原题“ionization energy”的符号应按其明确定义理解。冻结轨道是定理必要的近似。
\end{solution}
\originalsection{作业8（2004年11月16日截止）}
\begin{exercise}[原题1(a)--(c)]
(a)均匀正背景中的电子气，求HF本征值与基态能。(b)LDA中$E_{xc}=\int n\epsilon_{xc}(n)\dd r$，$\mu_{xc}=\dd(n\epsilon_{xc})/\dd n$，证明外势$v(r)$中的总能为$\sum_i\epsilon_i-E_H+\int n(\epsilon_{xc}-\mu_{xc})\dd r$。(c)把$\epsilon_{xc}$只取(a)的交换能，求均匀体系KS能级，并比较HF的费米速度、带宽和基态能；背景自能也须计入。
\end{exercise}
\begin{solution}
(a)第\ref{ch:interaction}章给$\epsilon_k=\hbar^2k^2/(2m)+\Sigma_x(k)$，$E/N=3E_F/5-3e_*^2k_F/(4\pi)$。正--正、正--负和电子Hartree长程项相互抵消，若漏背景自能则会剩无限体积发散。

(b)将KS方程按占据轨道求和得$\sum_i\epsilon_i=T_s+\int vn+2E_H+\int n\mu_{xc}$；与$E=T_s+\int vn+E_H+\int n\epsilon_{xc}$比较即得公式。

(c)$\epsilon_x=-3e_*^2k_F/(4\pi)$且$k_F\propto n^{1/3}$，所以$\mu_x=4\epsilon_x/3=-e_*^2k_F/\pi$。均匀KS能$\epsilon_k^{KS}=\hbar^2k^2/(2m)-e_*^2k_F/\pi$，交换是常数，速度为裸$\hbar k_F/m$，带宽$E_F$。HF自能在0与$k_F$间上升$e_*^2k_F/\pi$，故带宽$E_F+e_*^2k_F/\pi$，费米速度有对数发散。KS按(b)计算总能仍与HF交换总能相同：辅助能级差不等于总能泛函差。
\end{solution}
原题2要求阅读Marder附录C学习二次量子化，没有独立计算题。本册第\ref{ch:interaction}章提供所需先修，阅读指针不计为已取得的书本题。
\originalsection{作业9（2004年11月23日截止）}
\begin{exercise}[原题1--2]
两站Hubbard模型$H=t\sum_\sigma(c_{1\sigma}^\dagger c_{2\sigma}+\mathrm{h.c.})+U\sum_{i=1,2}n_{i\uparrow}n_{i\downarrow}$。(1)用创建算符写1、2、3、4电子Hilbert基及矩阵，给1、3、4电子本征能与态。(2)精确求2电子本征能与简并，并在$U\gg|t|$与自旋$1/2$的$J\vct S_1\cdot\vct S_2$比较，求$J$。
\end{exercise}
\begin{solution}
固定单轨道顺序$(1\uparrow,1\downarrow,2\uparrow,2\downarrow)$，各粒子数基为从四个$c_\alpha^\dagger$中选$N$个按顺序乘真空，维数依次$4,6,4,1$。矩阵元由$tc_\alpha^\dagger c_\beta$按反对易换序得，源码另以此规则实际构矩阵复核。

$N=1$每自旋矩阵$\left(\begin{smallmatrix}0&t\\t&0\end{smallmatrix}\right)$，态$(c_{1\sigma}^\dagger\pm c_{2\sigma}^\dagger)\ket0/\sqrt2$，能量$\pm t$。$N=3$按一个空穴基等价于对角$U$的两个$2\times2$块，能量$U\pm t$各二重，本征矢是两种相同剩余自旋配置的双占据位置叠加，相位由固定基顺序决定；严格向量与完整矩阵随计算证据给出。$N=4$唯一全占据态，跃迁因Pauli原理阻塞，能量$2U$。

$N=2$三重态$T_+=c_{1\uparrow}^\dagger c_{2\uparrow}^\dagger\ket0,T_-=c_{1\downarrow}^\dagger c_{2\downarrow}^\dagger\ket0,T_0=(c_{1\uparrow}^\dagger c_{2\downarrow}^\dagger+c_{1\downarrow}^\dagger c_{2\uparrow}^\dagger)\ket0/\sqrt2$皆能量0。单重态$S=(c_{1\uparrow}^\dagger c_{2\downarrow}^\dagger-c_{1\downarrow}^\dagger c_{2\uparrow}^\dagger)\ket0/\sqrt2$与$D_+=(c_{1\uparrow}^\dagger c_{1\downarrow}^\dagger+c_{2\uparrow}^\dagger c_{2\downarrow}^\dagger)\ket0/\sqrt2$构成矩阵$\left(\begin{smallmatrix}0&2t\\2t&U\end{smallmatrix}\right)$，本征根$E_\pm=(U\pm\sqrt{U^2+16t^2})/2$，态归一化为$S+(E_\pm/2t)D_+$。$D_-=(D_1-D_2)/\sqrt2$能量$U$。$t=0$用连续极限或直接原基处理。

强耦合$E_-=-4t^2/U+O(t^4/U^3)$，三重态仍0，匹配$J(\vct S_1\cdot\vct S_2-1/4)$，故$J=4t^2/U$。比较普通Heisenberg时需加常数$-J/4$，否则本征能绝对值不会相同。
\end{solution}
\originalsection{作业10（2004年11月30日截止）}
\begin{exercise}[原题1(a)--(b)]
零温费米能上方$E$的电子：(a)证明低能声子散射单粒子宽度量级$\hbar/\tau\sim E^3/(\hbar\omega_D)^2$；(b)输运时间在跃迁概率中乘$1-\cos\theta$，证明$\hbar/\tau_t\sim E^5/(\hbar\omega_D)^4$。只求幂律及自然尺度，材料系数可省略。
\end{exercise}
\begin{solution}
声学声子$\omega_q=c_sq$且$|g_q|^2\propto q$。三维相空间$q^2\dd q$，能量守恒角积分给$1/(v_Fq)$，零温Pauli因子限制$\hbar c_sq<E$，所以率$\propto\int_0^{E/\hbar c_s}q^2\dd q\propto E^3$。低温大角输运由小声子动量造成的电子转角$\theta\sim q/k_F$，故多$\int q^2$权重，率$\propto E^5$。以Debye能归一后量纲给题示尺度，但系数仍含耦合及$k_F/q_D$比；替换$E\to\kb T$只给有限温幂律，不给精确系数。该推导要求三维、声学形变耦合与$E\ll\hbar\omega_D$。
\end{solution}
\originalsection{作业11（2004年12月9日截止）}
\begin{exercise}[原题1(a)--(c)]
约化BCS哈密顿量$H=\sum_{k\sigma}\epsilon_kc_{k\sigma}^\dagger c_{k\sigma}-V_0\sum_{kk'}c_{k\uparrow}^\dagger c_{-k\downarrow}^\dagger c_{-k'\downarrow}c_{k'\uparrow}$与BCS态。(a)求$W=\avg{H-\mu N}$。(b)在$u_k^2+|v_k|^2=1$下变分，令$u=\cos\theta,v=\sin\theta$，求$\tan2\theta=\Delta/\xi$与零温隙方程。(c)从有限温方程推低温$\Delta(T)=\Delta_0[1-a(T)\e^{-\Delta_0/\kb T}]$，证明$a(T)\propto\sqrt T$。
\end{exercise}
\begin{solution}
(a)$\avg{n_{k\sigma}}=|v_k|^2$，配对平均$u_k^*v_k$，得$W=2\sum_k\xi_k|v_k|^2-V_0|\sum_ku_k^*v_k|^2$（约化模型的热力学极限，重复$k$自项的$O(1/V)$约定不改隙方程）。(b)统一相位令$\Delta$实，$\partial_{\theta_k}W=2\xi_k\sin2\theta_k-2\Delta\cos2\theta_k=0$，所以$\sin2\theta=\Delta/\sqrt{\xi^2+\Delta^2}$，$\Delta=V_0\sum\Delta/(2E_k)$。近费米面每自旋态密度常数给$1=N_\uparrow(0)V_0\int_0^{\hbar\omega_D}\dd\xi/\sqrt{\xi^2+\Delta^2}$。

(c)从有限温积分扣除零温积分，弱耦合下$\log(\Delta_0/\Delta)=2\int_0^\infty f(\sqrt{\xi^2+\Delta^2})\dd\xi/\sqrt{\xi^2+\Delta^2}$。低温以$E\simeq\Delta_0+\xi^2/(2\Delta_0)$估计高斯积分，右侧$\simeq\sqrt{2\pi\kb T/\Delta_0}\e^{-\Delta_0/\kb T}$，故
$\Delta_0-\Delta\simeq\sqrt{2\pi\Delta_0\kb T}\e^{-\Delta_0/\kb T}$。这里$a(T)=\sqrt{2\pi\kb T/\Delta_0}$，不是与温度无关常数；原题也说明其含$\sqrt T$。
\end{solution}
